%%% BEGIN fig:coll-photoproduction
\begin{figure}[htbp]
\centering
\begin{tikzpicture}[fd,
  fb lab/.style={font=\small},
  fb note/.style={font=\footnotesize,inner sep=1.5pt},
  fb hadron/.style={draw=black!70,line width=1.5pt,line cap=butt},
  fb hb/.style={hard blob,circle,minimum size=0.84cm,inner sep=0pt}]
  % ======================= (a) direct photoproduction =======================
  \begin{scope}
    \node[fb hb] (H) at (0,0) {$H$};
    \node[jet blob,ellipse,minimum width=1.25cm,minimum height=0.66cm,inner sep=0pt] (A) at (-1.55,-1.35) {};
    \node[font=\footnotesize] at (A) {$f_{i/p}$};
    \draw[photon] (-2.75,1.35) -- (H.145);
    \draw[fb hadron] (-3.05,-1.35) -- (A.west);
    \foreach \d in {-0.09,0.09} \draw[plain] ($(A.east)+(0,\d)$) -- ++(3.1,0);
    \draw[quark] (A.45) -- (H.235);
    \draw[quark] (H.25) -- (1.9,0.95);
    \draw[quark] (H.335) -- (1.9,-0.55);
    \node[fb lab,anchor=east] at (-2.77,1.35) {$\gamma$};
    \node[fb lab,anchor=east] at (-3.07,-1.35) {$p$};
    \node[fb note,anchor=west] at (2.2,-1.35) {$X_p$};
    \node[fb note,anchor=west,align=left] at (1.95,0.2) {高 $p_T$\\喷注};
    \path[mom={$\xi P$},mom sep=0.26cm,mom length=0.45cm] (A.45) -- (H.235);
    \node[font=\small,anchor=north,align=center] at (-0.35,-1.85) {(a)\ 直接光子：单强子过程\\{\color{RoyalBlue!80!black}全阶论证 [A]}};
  \end{scope}
  % ======================= (b) resolved photoproduction =======================
  \begin{scope}[xshift=8.0cm]
    \node[fb hb] (H) at (0,0) {$H$};
    \node[jet blob,ellipse,minimum width=1.25cm,minimum height=0.66cm,inner sep=0pt] (A) at (-1.55,-1.35) {};
    \node[jetB blob,ellipse,minimum width=1.25cm,minimum height=0.66cm,inner sep=0pt] (G) at (-1.55,1.35) {};
    \node[font=\footnotesize] at (A) {$f_{i/p}$};
    \node[font=\footnotesize] at (G) {$f_{j/\gamma}$};
    \draw[photon] (-3.05,1.35) -- (G.west);
    \draw[fb hadron] (-3.05,-1.35) -- (A.west);
    \foreach \d in {-0.09,0.09} {
      \draw[plain] ($(A.east)+(0,\d)$) -- ++(3.1,0);
      \draw[plain] ($(G.east)+(0,\d)$) -- ++(3.1,0);}
    \draw[quark] (A.45) -- (H.235);
    \draw[quark] (G.-45) -- (H.125);
    \draw[quark] (H.25) -- (1.75,0.8);
    \draw[quark] (H.335) -- (1.75,-0.8);
    \draw[glauber] (2.0,1.24) -- (2.0,-1.24);
    \node[fb lab,anchor=east] at (-3.07,1.35) {$\gamma$};
    \node[fb lab,anchor=east] at (-3.07,-1.35) {$p$};
    \node[fb note,anchor=west] at (2.2,-1.35) {$X_p$};
    \node[fb note,anchor=west] at (2.2,1.35) {$X_\gamma$};
    \node[fb note,anchor=west,align=left,text=red!55!black] at (2.12,0) {旁观者间的\\Glauber 交换};
    \path[mom={$\xi P$},mom sep=0.26cm,mom length=0.45cm] (A.45) -- (H.235);
    \path[mom'={$x_\gamma q$},mom sep=0.26cm,mom length=0.45cm] (G.-45) -- (H.125);
    \node[font=\small,anchor=north,align=center] at (-0.35,-1.85) {(b)\ 分辨光子：类强子—强子过程\\{\color{Plum!80!black}猜想 [C]}};
  \end{scope}
\end{tikzpicture}
\caption{光生过程 $\gamma p\to\text{jets}+X$（振幅示意）。(a)~直接光子：光子直接参与硬散射 $H$，只有质子一侧有部分子分布 $f_{i/p}$ 与旁观者 $X_p$，与 DIS 同为单强子过程，因子化为全阶论证 [A]（\cref{thm:coll-direct-photo}）；
(b)~分辨光子：光子先“强子化”为由光子部分子分布 $f_{j/\gamma}$ 描述的部分子云，提供进入硬散射的部分子（动量份额 $x_\gamma$），光子一侧也有旁观者 $X_\gamma$，
等效于末态有色的强子—强子过程，旁观者之间可交换 Glauber 胶子（暗红色），因子化为猜想 [C]（\cref{conj:coll-resolved-photo}）。}
\label{fig:coll-photoproduction}
\end{figure}
%%% END fig:coll-photoproduction
